\documentclass[12pt,a4paper]{article}

% ── Paquetes ──────────────────────────────────────────────────────────────────
\usepackage[utf8]{inputenc}
\usepackage[T1]{fontenc}
\usepackage[spanish]{babel}
\usepackage{geometry}
\usepackage{graphicx}
\usepackage{hyperref}
\usepackage{listings}
\usepackage{xcolor}
\usepackage{amsmath}
\usepackage{amssymb}
\usepackage{booktabs}
\usepackage{tabularx}
\usepackage{float}
\usepackage{enumitem}
\usepackage{fancyhdr}
\usepackage{titlesec}
% \usepackage{microtype}  % desactivado: incompatible con fuentes bitmap en este sistema
\usepackage{parskip}
\usepackage{cite}
\usepackage{caption}
\usepackage{subcaption}
\usepackage{mdframed}
\usepackage{array}
\usepackage{multirow}
\usepackage{tikz}
\usetikzlibrary{arrows.meta,automata,positioning,shapes.geometric}

% ── Geometría ────────────────────────────────────────────────────────────────
\geometry{
  top=2.5cm, bottom=2.5cm,
  left=3cm,  right=2.5cm,
  headheight=14pt
}

% ── Colores ──────────────────────────────────────────────────────────────────
\definecolor{codeblue}{RGB}{0,73,144}
\definecolor{codegray}{RGB}{100,100,100}
\definecolor{codegreen}{RGB}{0,128,0}
\definecolor{codepurple}{RGB}{148,0,211}
\definecolor{backcolor}{RGB}{248,248,248}
\definecolor{hulkred}{RGB}{170,0,0}
\definecolor{rustbrown}{RGB}{183,65,14}

% ── Estilo de listados ───────────────────────────────────────────────────────
\lstdefinestyle{hulk}{
  backgroundcolor=\color{backcolor},
  basicstyle=\ttfamily\footnotesize,
  keywordstyle=\color{codeblue}\bfseries,
  stringstyle=\color{codegreen},
  commentstyle=\color{codegray}\itshape,
  numberstyle=\tiny\color{codegray},
  breaklines=true,
  breakatwhitespace=true,
  captionpos=b,
  frame=single,
  framerule=0.5pt,
  numbers=left,
  numbersep=8pt,
  showstringspaces=false,
  tabsize=4,
  morekeywords={let,in,if,elif,else,while,for,function,type,new,self,inherits,
    protocol,extends,is,as,true,false,Number,Boolean,String,Object,Null,Vector},
}

\lstdefinestyle{rust}{
  backgroundcolor=\color{backcolor},
  basicstyle=\ttfamily\footnotesize,
  keywordstyle=\color{rustbrown}\bfseries,
  stringstyle=\color{codegreen},
  commentstyle=\color{codegray}\itshape,
  numberstyle=\tiny\color{codegray},
  breaklines=true,
  frame=single,
  framerule=0.5pt,
  numbers=left,
  numbersep=8pt,
  showstringspaces=false,
  tabsize=4,
  morekeywords={fn,let,mut,pub,struct,enum,impl,match,Ok,Err,Some,None,
    use,mod,crate,super,self,return,if,else,while,for,in,loop},
}

\lstdefinestyle{llvmir}{
  backgroundcolor=\color{backcolor},
  basicstyle=\ttfamily\footnotesize,
  keywordstyle=\color{codepurple}\bfseries,
  commentstyle=\color{codegray}\itshape,
  frame=single,
  framerule=0.5pt,
  numbers=left,
  numbersep=8pt,
  breaklines=true,
  morekeywords={define,declare,alloca,store,load,call,ret,br,add,sub,mul,
    fdiv,fadd,fmul,fsub,fptoui,fptosi,uitofp,getelementptr,icmp,fcmp,phi,
    unreachable,i1,i8,i32,i64,double,ptr,void,global,constant},
}

\lstset{style=hulk}

% ── Encabezados y pie de página ──────────────────────────────────────────────
\pagestyle{fancy}
\fancyhf{}
\fancyhead[L]{\small\textit{Compilador HULK --- Reporte de Proyecto}}
\fancyhead[R]{\small\textit{Facultad de Matematica y Computacion, UH}}
\fancyfoot[C]{\thepage}
\renewcommand{\headrulewidth}{0.4pt}

% ── Hiperlinks ───────────────────────────────────────────────────────────────
\hypersetup{
  colorlinks=true,
  linkcolor=codeblue,
  citecolor=codeblue,
  urlcolor=codeblue,
  pdftitle={Compilador HULK -- Reporte de Proyecto},
  pdfauthor={Equipo HULK},
}

% ── Titulo ───────────────────────────────────────────────────────────────────
\title{
  \vspace{-1cm}
  \Huge\textbf{Compilador HULK}\\[0.5em]
  \large Reporte de Proyecto Final\\[0.3em]
  \large Asignatura: Compilacion / Lenguajes de Programacion
}
\author{
  \textbf{Facultad de Matematica y Computacion}\\
  Universidad de La Habana\\[1em]
  Ciencias de la Computacion --- Tercer Ano
}
\date{\today}

% ─────────────────────────────────────────────────────────────────────────────
\begin{document}
% ─────────────────────────────────────────────────────────────────────────────

\maketitle
\thispagestyle{empty}
\vspace{2cm}

\begin{abstract}
\noindent
Este documento describe el diseno, la arquitectura y la implementacion de un compilador
completo para el lenguaje \textbf{HULK} (Havana University Language for Kids), desarrollado
enteramente en el lenguaje de programacion \textbf{Rust} con backend LLVM~17 mediante la
biblioteca \textit{inkwell}. El compilador cubre todos los requerimientos obligatorios de la
especificacion HULK (secciones~1 a~16.8, incluyendo verificacion de tipos) y extiende el
lenguaje con un \textbf{sistema de vectores con expresiones generadoras}, una extension
sintactica y semantica con base en la teoria de comprensiones de listas presentes en
lenguajes como Haskell y Python. El documento incluye una descripcion detallada de cada
fase del compilador, las decisiones de diseno, comparativas con otros lenguajes para las
extensiones propuestas, y una evaluacion de la cobertura de la suite de pruebas oficial.
\end{abstract}

\vfill
\begin{mdframed}
\small
\textbf{Tecnologias utilizadas:} Rust 1.77+, LLVM 17, inkwell 0.4, LALR(1) artesanal,
NFA por construccion de Thompson.\\
\textbf{Lineas de codigo Rust:} $\approx$8\,000 LOC fuente + $\approx$2\,000 LOC de pruebas.\\
\textbf{Tests automatizados:} 384 pruebas en \texttt{cargo test} + suite oficial del profesor.
\end{mdframed}

\newpage
\tableofcontents
\newpage

% =============================================================================
\section{Introduccion}
% =============================================================================

HULK es un lenguaje de programacion disenado con fines educativos en la Facultad de
Matematica y Computacion de la Universidad de La Habana. Es un lenguaje orientado a
expresiones, estaticamente tipado, con soporte para programacion orientada a objetos
(herencia simple, polimorfismo, protocolos) y construcciones funcionales (funciones de
primera clase, expresiones generadoras).

\subsection{Objetivos del proyecto}

El objetivo principal es implementar un compilador de produccion para HULK que:
\begin{enumerate}
  \item Cubra \textbf{todos} los features obligatorios de la especificacion (hasta la
    seccion~16.8, verificacion de tipos incluida).
  \item Genere \textbf{codigo nativo} x86-64 mediante LLVM~17 para obtener ejecutables
    eficientes que no requieran un interprete en tiempo de ejecucion.
  \item Implemente una \textbf{extension propia} al lenguaje que modifique tanto la
    sintaxis como la semantica, con fundamentacion teorica y comparativa respecto a
    otros lenguajes.
  \item Sea completamente desarrollado en \textbf{Rust}, aprovechando las garantias de
    seguridad de memoria del lenguaje anfitrion.
\end{enumerate}

\subsection{Estructura del compilador}

El compilador sigue el diseno clasico de compiladores en fases encadenadas:

\begin{figure}[H]
\centering
\begin{tikzpicture}[node distance=1.1cm, auto,
  box/.style={rectangle, draw=codeblue, fill=blue!5, rounded corners,
    minimum width=2.8cm, minimum height=0.75cm, font=\small},
  arr/.style={-Stealth, thick, color=codeblue}]
  \node[box] (src)  {Fuente .hulk};
  \node[box] (lex)  [below=of src]  {Analisis lexico};
  \node[box] (par)  [below=of lex]  {Analisis sintactico};
  \node[box] (sem)  [below=of par]  {Analisis semantico};
  \node[box] (cg)   [below=of sem]  {Generacion de codigo};
  \node[box] (opt)  [below=of cg]   {Optimizacion LLVM};
  \node[box] (bin)  [below=of opt]  {Binario nativo};
  \foreach \a/\b in {src/lex, lex/par, par/sem, sem/cg, cg/opt, opt/bin}
    \draw[arr] (\a) -- (\b);
  \node[right=0.5cm of lex, font=\scriptsize, text=codegray]  {NFA Thompson + subset};
  \node[right=0.5cm of par, font=\scriptsize, text=codegray]  {LALR(1) artesanal};
  \node[right=0.5cm of sem, font=\scriptsize, text=codegray]  {Dos pasadas};
  \node[right=0.5cm of cg,  font=\scriptsize, text=codegray]  {LLVM IR via inkwell};
  \node[right=0.5cm of opt, font=\scriptsize, text=codegray]  {mem2reg, instcombine};
\end{tikzpicture}
\caption{Pipeline de compilacion de HULK.}
\end{figure}

Cada fase es un modulo Rust independiente con interfaces bien definidas, lo que
facilita el desarrollo incremental, las pruebas unitarias y el mantenimiento.

\subsection{Eleccion del lenguaje: Rust}

La especificacion permite C, C++ o Rust. Se eligio Rust por tres razones principales:

\begin{itemize}
  \item \textbf{Seguridad de memoria sin GC}: el compilador maneja estructuras de datos
    complejas con muchas referencias cruzadas (tablas de simbolos, jerarquias de tipos,
    tabla de parseo LALR). Rust garantiza ausencia de dangling pointers y data races en
    tiempo de compilacion, eliminando una clase entera de bugs dificiles de depurar.
  \item \textbf{Bindings seguros a LLVM}: la biblioteca \textit{inkwell} proporciona una
    API de alto nivel que envuelve la C API de LLVM con tipos Rust seguros, evitando
    los peligros del FFI crudo de C++.
  \item \textbf{Sistemas de tipos expresivos}: los enumerados con datos asociados
    (\texttt{enum}) y el \textit{pattern matching} exhaustivo hacen que representar y
    manipular el AST sea natural y seguro.
\end{itemize}

% =============================================================================
\section{Analisis Lexico}
% =============================================================================

\subsection{Arquitectura: Lexer basado en NFA}

El lexer se implementa usando la \textbf{construccion de Thompson}~\cite{thompson1968} para
convertir expresiones regulares en Automatas Finitos No-deterministicos (NFA). A diferencia
del enfoque clasico que convierte el NFA a DFA antes de usarlo, este lexer ejecuta el NFA
directamente usando \textit{construccion de subconjuntos en tiempo de ejecucion}.

La arquitectura tiene tres capas:

\begin{enumerate}
  \item \textbf{Capa de regex}: un mini-parser de expresiones regulares que produce un AST
    de regex (\texttt{RegexNode}). Se soportan concatenacion, alternacion (\texttt{|}),
    clausura de Kleene (\texttt{*}), clausura positiva (\texttt{+}), opcionalidad
    (\texttt{?}), rangos de caracteres (\texttt{[a-z]}), y el punto (\texttt{.}).
  \item \textbf{Construccion de Thompson}: \texttt{thompson.rs} transforma el AST de regex
    en un NFA con transiciones epsilon segun las reglas estandar: cada operador produce
    un fragmento NFA con exactamente un estado inicial y uno de aceptacion.
  \item \textbf{MasterNFA}: \texttt{master\_nfa.rs} combina todos los NFA individuales
    (uno por tipo de token) en un unico NFA mediante transiciones epsilon desde un estado
    inicial compartido. Durante la tokenizacion, el motor ejecuta la construccion de
    subconjuntos sobre este NFA para encontrar el token mas largo.
\end{enumerate}

\begin{figure}[H]
\centering
\begin{tikzpicture}[auto, node distance=1.5cm,
  state/.style={circle, draw=codeblue, fill=blue!5, minimum size=0.7cm, font=\footnotesize},
  acc/.style={circle, draw=codeblue, fill=blue!20, minimum size=0.7cm,
    font=\footnotesize, double},
  arr/.style={-Stealth, thick}]
  \node[state] (s0) {$q_0$};
  \node[state] (s1) [right=of s0] {$q_1$};
  \node[state] (s2) [right=of s1] {$q_2$};
  \node[state] (s3) [right=of s2] {$q_3$};
  \node[acc]   (s4) [right=of s3] {$q_4$};
  \draw[arr] (s0) -- node{\texttt{i}} (s1);
  \draw[arr] (s1) -- node{\texttt{f}} (s2);
  \draw[arr] (s2) -- node{$\varepsilon$} (s3);
  \draw[arr] (s0) edge[bend right=25] node[below]{\texttt{i}} (s3);
  \draw[arr] (s3) -- node{\texttt{d}} (s4);
  \node[font=\scriptsize, text=codegray] at (2.5,-1.2) {NFA para \texttt{if}
    $\cup$ \texttt{id}: la prioridad decide};
\end{tikzpicture}
\caption{Ejemplo simplificado: NFA combinado para \texttt{if} e identificadores.}
\end{figure}

La resolucion de ambiguedad (p.ej. \texttt{if} vs identificador) se hace por
\textit{prioridad}: cada patron tiene un numero de prioridad; cuando multiples NFA
aceptan la misma cadena, se escoge el de mayor prioridad (keywords antes que identificadores).
En caso de igual prioridad, gana el que haya reconocido la cadena mas larga (regla del
maximo prefijo, \textit{maximal munch}).

\subsection{Conjunto de tokens}

La tabla siguiente resume los grupos de tokens reconocidos:

\begin{table}[H]
\centering
\begin{tabular}{lll}
\toprule
\textbf{Categoria} & \textbf{Ejemplos} & \textbf{Descripcion} \\
\midrule
Keywords    & \texttt{let}, \texttt{in}, \texttt{if}, \texttt{while},
              \texttt{for}, \texttt{type}, \texttt{protocol} & Palabras reservadas \\
Operadores  & \texttt{+}, \texttt{-}, \texttt{*}, \texttt{/}, \texttt{\%},
              \texttt{\^{}}, \texttt{@}, \texttt{@@}, \texttt{:=} & Aritm., concat., asig. \\
Comparacion & \texttt{==}, \texttt{!=}, \texttt{<}, \texttt{>},
              \texttt{<=}, \texttt{>=} & Relacionales \\
Logicos     & \texttt{\&\&}, \texttt{||}, \texttt{!} & Booleanos \\
OOP         & \texttt{new}, \texttt{self}, \texttt{inherits}, \texttt{is},
              \texttt{as}, \texttt{base} & Orientacion a objetos \\
Literales   & \texttt{42}, \texttt{3.14}, \texttt{"hola"}, \texttt{true} & Numeros, cadenas, bool \\
Delimitadores & \texttt{(}, \texttt{)}, \texttt{\{}, \texttt{\}},
                \texttt{[}, \texttt{]}, \texttt{,}, \texttt{;}, \texttt{:} & Estructura \\
\bottomrule
\end{tabular}
\caption{Categorias de tokens del lexer HULK.}
\end{table}

\subsection{Procesamiento de secuencias de escape}

Los literales de cadena soportan las secuencias de escape estandar. El procesamiento
se realiza en la fase de acciones semanticas del parser, cuando se construye el nodo
\texttt{Literal::String} del AST:

\begin{lstlisting}[style=rust, caption={Procesamiento de escapes en semantic\_actions.rs}]
let content = {
    let raw = &lexeme[1..lexeme.len() - 1]; // quitar comillas
    let mut out = String::with_capacity(raw.len());
    let mut chars = raw.chars().peekable();
    while let Some(c) = chars.next() {
        if c == '\\' {
            match chars.next() {
                Some('n')  => out.push('\n'),
                Some('t')  => out.push('\t'),
                Some('"')  => out.push('"'),
                Some('\\') => out.push('\\'),
                Some(o)    => { out.push('\\'); out.push(o); }
                None       => out.push('\\'),
            }
        } else { out.push(c); }
    }
    out
};
\end{lstlisting}

Se procesan en el AST (no en el lexer) porque \texttt{build\_global\_string\_ptr} de LLVM
recibe el valor ya interpretado; si se procesaran en el lexer, la informacion de posicion
de los tokens quedaria desincronizada con el texto fuente.

\subsection{Reporte de errores lexicos}

Los caracteres no reconocidos producen un token \texttt{ERROR} en lugar de abortar la
compilacion, permitiendo reportar multiples errores lexicos en una sola pasada. El formato
de error sigue el contrato de la interfaz oficial:

\begin{verbatim}
! (line,col) LEXICAL: unexpected character '$'
\end{verbatim}

% =============================================================================
\section{Analisis Sintactico}
% =============================================================================

\subsection{Parser LALR(1) implementado desde cero}

En lugar de usar un generador de parsers externo (como YACC, Bison, LALRPOP o ANTLR),
el compilador implementa el algoritmo LALR(1) completo en Rust. Esta decision se
detalla en la seccion~\ref{sec:decisions}.

El modulo de parseo esta organizado en cuatro capas:

\begin{enumerate}
  \item \textbf{Definicion de la gramatica} (\texttt{grammar/}): la gramatica de HULK
    se define como estructura de datos Rust (\texttt{Vec<Production>}), no como un DSL.
    Esto hace que la gramatica sea inspeccionable en tiempo de ejecucion, lo que
    facilita la depuracion.
  \item \textbf{Calculo de FIRST y FOLLOW} (\texttt{lalr/first\_follow.rs}):
    algoritmo de punto fijo estandar para calcular los conjuntos
    $\text{FIRST}(\alpha)$ y $\text{FOLLOW}(A)$.
  \item \textbf{Construccion de la tabla LALR} (\texttt{lalr/table\_builder.rs}):
    coleccion canonica de items LR(0), calculo de lookaheads LALR por propagacion,
    construccion de las tablas ACTION/GOTO.
  \item \textbf{Motor del parser} (\texttt{engine/}): algoritmo shift-reduce con
    acciones semanticas invocadas en cada reduccion para construir el AST.
\end{enumerate}

\subsection{Jerarquia de la gramatica y precedencia de operadores}

La especificacion original de HULK deja la precedencia de operadores implicita. La
gramatica implementada resuelve la ambiguedad mediante la \textit{jerarquia de
no-terminales}, codificando la precedencia directamente en la estructura de la gramatica:

\begin{lstlisting}[style={}, caption={Jerarquia de precedencias en la gramatica HULK},
  numbers=none]
expr           -> assignment
assignment     -> or_expr (':=' assignment)?
or_expr        -> and_expr ('||' and_expr)*
and_expr       -> not_expr ('&&' not_expr)*
not_expr       -> '!' not_expr | equality_expr
equality_expr  -> comparison (('=='|'!=') comparison)*
comparison     -> concat (('<'|'>'|'<='|'>=') concat)*
concat         -> additive (('@'|'@@') additive)*
additive       -> term (('+' | '-') term)*
term           -> power (('*' | '/' | '%') power)*
power          -> unary ('^' power)?       -- asociatividad derecha
unary          -> '-' unary | postfix
postfix        -> atom ('.' id '(' args ')' | '.' id | '[' expr ']')*
atom           -> literal | id | '(' expr ')' | ...
\end{lstlisting}

Esta codificacion garantiza la precedencia estandar matematica (potencia > producto >
suma > comparacion > logicos) sin necesidad de reglas de desambiguacion adicionales.
La potencia tiene asociatividad derecha ($2^{3^2} = 2^9$, no $(2^3)^2$), lo que se
implementa con la regla recursiva a la derecha.

\subsection{Arbol de Sintaxis Abstracta (AST)}

El AST se define en \texttt{src/parser/ast/} usando enumerados Rust con datos
asociados. El nodo central es \texttt{Expr}, que contiene un \texttt{ExprKind} y
un \texttt{Span} (linea y columna) para todos los mensajes de error posteriores:

\begin{lstlisting}[style=rust, caption={Estructura principal del AST}]
pub struct Expr {
    pub kind: ExprKind,
    pub id:   u32,    // identificador unico para tabla de tipos
    pub span: Span,   // posicion en el fuente
}

pub enum ExprKind {
    Literal(Literal),
    Identifier { name: String },
    Binary(BinaryExpr),
    Unary(UnaryExpr),
    Let(LetExpr),
    If(IfExpr),
    While(WhileExpr),
    For(ForExpr),
    Block(BlockExpr),
    Call(CallExpr),
    MethodCall(MethodCallExpr),
    Access(AccessExpr),
    New(NewExpr),
    Is { expr: Box<Expr>, type_name: TypeName },
    As { expr: Box<Expr>, type_name: TypeName },
    Vector(VectorExpr),  // <-- extension: vectores y generadores
    Index(IndexExpr),    // <-- extension: indexacion v[i]
    Assign(AssignExpr),
    Postfix(PostfixExpr),
    Base,
}
\end{lstlisting}

Cada nodo lleva un identificador unico (\texttt{id}) que sirve como clave en la tabla
\texttt{expr\_types: HashMap<u32, HulkType>} que el analizador semantico construye.
El generador de codigo consulta esta tabla para conocer el tipo inferido de cada
subexpresion sin tener que reinferir.

% =============================================================================
\section{Analisis Semantico}
% =============================================================================

\subsection{Verificacion de tipos en dos pasadas}

El analizador semantico opera en \textbf{dos pasadas} sobre la lista de declaraciones
del programa:

\textbf{Pasada 1 --- Recoleccion de declaraciones:} se registran todas las firmas de
funciones, los miembros de tipos y los cuerpos de protocolos en la estructura
\texttt{TypeHierarchy}. Los miembros sin anotacion de tipo reciben el marcador
\texttt{HulkType::Unknown}. Esta pasada es necesaria para soportar funciones y tipos
mutuamente recursivos sin restricciones de orden.

\textbf{Pasada 2 --- Verificacion de cuerpos:} cada cuerpo de funcion, inicializador
de atributo y cuerpo de metodo se verifica. El tipo inferido se compara con la
anotacion declarada si la hay. Los tipos inferidos se \textbf{escriben de vuelta} a la
\texttt{TypeHierarchy}, reemplazando los marcadores \texttt{Unknown}. Esta escritura
es critica para el generador de codigo: sin ella, los atributos sin anotacion mantienen
tipo \texttt{Unknown} y el codegen no puede determinar el tipo LLVM correcto.

\begin{lstlisting}[style=rust, caption={Fragmento de la escritura de tipos inferidos}]
// Despues de inferir el tipo del campo en la pasada 2:
if let Some(type_info) = self.types.types.get_mut(type_name) {
    if let Some(attr) = type_info.attributes.get_mut(field_name) {
        if *attr == HulkType::Unknown {
            *attr = inferred_type.clone();
        }
    }
}
\end{lstlisting}

\subsection{Sistema de tipos}

El sistema de tipos es \textbf{nominal} con \textit{subtype polymorphism}:

\begin{itemize}
  \item \textbf{Tipos primitivos}: \texttt{Number} (IEEE 754 double), \texttt{Boolean}
    (i1), \texttt{String} (puntero a cadena null-terminated), \texttt{Null}.
  \item \textbf{Tipos definidos por el usuario}: herencia simple via \texttt{inherits}.
    Cada tipo ocupa exactamente un lugar en la jerarquia.
  \item \textbf{Protocolos}: tipos de interfaz estructural. Un tipo conforma a un
    protocolo si declara todos los metodos requeridos con firmas compatibles. A
    diferencia de las interfaces Java/C\#, los protocolos en HULK no requieren
    declaracion explicita de conformidad: la verificacion es estructural.
  \item \textbf{Vectores}: el tipo \texttt{[T]} representa un vector de elementos de
    tipo \texttt{T}. Es un tipo generico parametrico con un parametro de tipo.
  \item \textbf{HulkType::Unknown}: marcador para tipos no resueltos aun en la pasada~1.
  \item \textbf{HulkType::Never}: tipo de error que se propaga silenciosamente para
    evitar cascadas de errores falsos.
\end{itemize}

El predicado de \textit{conformidad} $A \leq B$ se define inductivamente:
$$A \leq A \qquad \frac{A \text{ hereda de } P \quad P \leq B}{A \leq B}
\qquad \frac{A \text{ implementa protocolo } P}{A \leq P}$$

\subsection{Funciones y constantes predefinidas}

El compilador pre-registra los siguientes elementos en la tabla de simbolos global:

\begin{table}[H]
\centering
\begin{tabular}{lll}
\toprule
\textbf{Nombre} & \textbf{Tipo} & \textbf{Descripcion} \\
\midrule
\texttt{print(x: Object)} & \texttt{Null} & Imprime cualquier valor a stdout \\
\texttt{sqrt(x: Number)} & \texttt{Number} & Raiz cuadrada \\
\texttt{sin(x: Number)} & \texttt{Number} & Seno (radianes) \\
\texttt{cos(x: Number)} & \texttt{Number} & Coseno (radianes) \\
\texttt{exp(x: Number)} & \texttt{Number} & Exponencial ($e^x$) \\
\texttt{log(b, x: Number)} & \texttt{Number} & Logaritmo en base $b$ \\
\texttt{rand(): Number} & \texttt{Number} & Aleatorio uniforme en $[0,1)$ \\
\texttt{range(s, e: Number)} & \texttt{Range} & Rango iterable $[s, e)$ \\
\texttt{PI} & \texttt{Number} & $\pi \approx 3.14159\ldots$ \\
\texttt{E}  & \texttt{Number} & $e \approx 2.71828\ldots$ \\
\bottomrule
\end{tabular}
\caption{Funciones y constantes predefinidas de HULK.}
\end{table}

\subsection{Reporte de errores semanticos}

Todos los errores semanticos se acumulan en un \texttt{Vec<SemanticError>} y se
reportan todos juntos al final del analisis, evitando que el usuario tenga que
corregir un error a la vez. Cada \texttt{SemanticError} lleva el \texttt{Span} de la
expresion o declaracion causante, produciendo mensajes del tipo:

\begin{verbatim}
! (12,5) SEMANTIC: type 'Cat' does not implement method 'speak' of protocol 'Animal'
! (18,3) SEMANTIC: cannot assign value of type 'String' to variable of type 'Number'
\end{verbatim}

% =============================================================================
\section{Generacion de Codigo}
% =============================================================================

\subsection{LLVM IR mediante inkwell}

El backend genera \textbf{Representacion Intermedia de LLVM} (LLVM IR) usando la
biblioteca \textit{inkwell}~\cite{inkwell}, que proporciona bindings seguros para la
API C de LLVM~17. El IR se optimiza y compila a codigo maquina nativo.

El estado de la generacion de codigo se centraliza en \texttt{CodegenContext}, que
agrupa:
\begin{itemize}
  \item \texttt{context}, \texttt{module}, \texttt{builder}: API LLVM central.
  \item \texttt{symbols}: pila de alcances (scope stack) para variables locales.
  \item \texttt{type\_registry}: layouts de structs LLVM y vtables por tipo.
  \item \texttt{type\_hierarchy}: jerarquia de tipos del analizador semantico.
  \item \texttt{expr\_types}: tipos anotados por el type checker, consultados en codegen.
  \item \texttt{current\_function}, \texttt{self\_ptr}: contexto de la funcion actual.
\end{itemize}

\subsection{Representacion de valores (CgValue)}

Dado que HULK tiene tipos heterogeneos, no existe un ``tipo universal'' de LLVM
para todos los valores. El enum \texttt{CgValue} transporta el valor LLVM junto con
su tipo HULK:

\begin{lstlisting}[style=rust, caption={Enum CgValue -- representacion de valores HULK en LLVM}]
pub enum CgValue<'ctx> {
    Number(FloatValue<'ctx>),    // LLVM: double
    Bool(IntValue<'ctx>),        // LLVM: i1
    Str(PointerValue<'ctx>),     // LLVM: ptr (C string null-terminated)
    Object(PointerValue<'ctx>),  // LLVM: ptr al struct del tipo
    Vector(PointerValue<'ctx>),  // LLVM: ptr al buffer de heap
    Null,                        // puntero nulo
    Void,                        // sin valor (efectos de lado)
}
\end{lstlisting}

Las coerciones entre variantes se realizan en el metodo \texttt{coerce\_arg}, que
inserta las instrucciones LLVM necesarias: por ejemplo, \texttt{Bool} a \texttt{double}
usa \texttt{uitofp}, y cualquier tipo primitivo a \texttt{ptr} para paso a funciones
polimorficas usa \texttt{alloca + store}.

\subsection{Representacion de objetos en memoria}

Todo tipo definido por el usuario se representa como un struct LLVM con la siguiente
distribucion de campos:

\begin{lstlisting}[style=llvmir, caption={Layout de struct para un tipo HULK con campos},
  numbers=none]
%MyType = type {
    i32,   ; campo 0: type_tag -- identifica el tipo en runtime
    ptr,   ; campo 1: vtable_ptr -- puntero a la vtable del tipo concreto
    double,; campo 2: primer campo del tipo (o del padre si hereda)
    ptr,   ; campo 3: segundo campo ...
    ...
}
\end{lstlisting}

La posicion de los campos en herencia sigue el orden \textit{padre primero}: todos los
campos del tipo padre aparecen antes de los campos propios, y esto se calcula
recursivamente. Esto garantiza compatibilidad de punteros: un puntero a un objeto hijo
puede tratarse como puntero al padre accediendo solo a los primeros campos.

\subsection{Vtables y despacho virtual}

El polimorfismo en HULK se implementa mediante \textbf{tablas de metodos virtuales
(vtables)}, una tecnica clasica usada por C++, Java y otros lenguajes compilados.

Para cada tipo se crea una vtable global que es un struct de punteros a funcion:

\begin{lstlisting}[style=llvmir, caption={Vtable de un tipo Dog que hereda de Animal}]
%VTable_Dog = type { ptr, ptr }  ; [speak, move]

@vtable_Dog = global %VTable_Dog {
    ptr @__hulk_method_Dog_speak,   ; override
    ptr @__hulk_method_Animal_move  ; heredado
}
\end{lstlisting}

El despacho de un metodo se realiza en tres pasos: (1) cargar el campo
\texttt{vtable\_ptr} del objeto, (2) calcular el slot del metodo en la vtable
(indice fijo conocido en tiempo de compilacion), (3) cargar el puntero a funcion
y llamarlo indirectamente:

\begin{lstlisting}[style=llvmir, caption={Despacho virtual de animal.speak()}]
  %vtable = load ptr, ptr %vtable_field
  %fn_slot = getelementptr %VTable_Animal, ptr %vtable, i32 0, i32 0
  %fn_ptr  = load ptr, ptr %fn_slot
  %result  = call double %fn_ptr(ptr %obj, ...)
\end{lstlisting}

\subsection{Type tags DFS y operador \texttt{is}}

El operador \texttt{is} (\texttt{x is Dog}) es una operacion frecuente en codigo OOP
con herencia. Implementarlo eficientemente requiere un invariante sobre la asignacion
de tags.

\textbf{Asignacion DFS pre-order:} durante la compilacion, los tipos reciben tags
enteros en orden de recorrido en profundidad de la jerarquia de herencia. El tipo padre
recibe su tag antes que sus hijos, por lo que todos los subtipos directos e indirectos
de un tipo $T$ forman un \textit{intervalo contiguo} $[\text{min\_tag}_T, \text{max\_tag}_T]$:

\begin{figure}[H]
\centering
\begin{tikzpicture}[level distance=1.2cm,
  every node/.style={rectangle, draw=codeblue, fill=blue!5,
    minimum width=1.8cm, minimum height=0.5cm, font=\small},
  edge from parent/.style={draw, -Stealth}]
  \node {Animal [1,4]}
    child { node {Dog [2,3]}
      child { node {Poodle [3,3]} }
    }
    child { node {Cat [4,4]} };
  \node[right=4cm, draw=none, fill=none, align=left, font=\small] at (0,0)
    {\texttt{poodle is Dog}: $2 \leq 3 \leq 3$ \checkmark\\
     \texttt{poodle is Animal}: $1 \leq 3 \leq 4$ \checkmark\\
     \texttt{cat is Dog}: $2 \leq 4 \leq 3$ \texttimes};
\end{tikzpicture}
\caption{Asignacion DFS de type tags. El test \texttt{is} es un range check O(1).}
\end{figure}

La instruccion LLVM generada para \texttt{x is Dog} es simplemente:

\begin{lstlisting}[style=llvmir, numbers=none]
%tag = load i32, ptr %x
%ge  = icmp uge i32 %tag, 2    ; min_tag(Dog)
%le  = icmp ule i32 %tag, 3    ; max_tag(Dog)
%res = and i1 %ge, %le
\end{lstlisting}

Esto es $O(1)$ en tiempo de ejecucion, sin consultar ninguna tabla ni recorrer la
jerarquia. Esta tecnica se utiliza en implementaciones de C++ con RTTI~\cite{bjarne1994}.

\subsection{Operador \texttt{as} -- downcast con verificacion en runtime}

El operador \texttt{x as T} realiza un \textit{downcast}: convierte un valor de tipo
base (p.ej. \texttt{Animal}) al tipo derivado \texttt{T} (p.ej. \texttt{Dog}).
En HULK esto es seguro porque si el objeto no es de tipo \texttt{T}, la ejecucion
se aborta con un mensaje de error.

El codigo generado primero verifica el tag, y si falla llama a \texttt{hulk\_type\_error}
(que imprime el mensaje y llama a \texttt{abort(3)}). Si el tag es valido, el puntero
original se usa directamente (no hay reinterpretacion de bits: el layout es compatible
por la organizacion padre-primero).

\subsection{Llamada a \texttt{base()}}

La palabra clave \texttt{base(args)} dentro de un metodo invoca la implementacion del
padre. A diferencia del despacho virtual (que cargaria la vtable y podria hacer
una llamada recursiva), \texttt{base()} es una \textbf{llamada directa estatica}:
el nombre de la funcion del padre se conoce en tiempo de compilacion.

El codegen busca recursivamente en la jerarquia hasta encontrar el tipo que implementa
el metodo y genera:
\begin{lstlisting}[style=llvmir, numbers=none]
%result = call double @__hulk_method_Animal_speak(ptr %self, ...)
\end{lstlisting}

\subsection{Pipeline de optimizacion}

El IR generado por el codegen es \textit{alloca-heavy}: todas las variables locales
se modelan como \texttt{alloca + store + load}. Antes de emitir codigo maquina se
aplica el pipeline de optimizacion del nuevo pass manager de LLVM~17:

\begin{table}[H]
\centering
\begin{tabular}{ll}
\toprule
\textbf{Pass} & \textbf{Efecto} \\
\midrule
\texttt{mem2reg} & Convierte \texttt{alloca/store/load} a registros SSA; base de todo lo demas \\
\texttt{instcombine} & Plegado de constantes, simplificacion algebraica \\
\texttt{reassociate} & Reordena expresiones para maxima oportunidad de plegado \\
\texttt{simplifycfg} & Elimina bloques basicos muertos, simplifica el CFG \\
\bottomrule
\end{tabular}
\caption{Passes de optimizacion LLVM aplicados.}
\end{table}

% =============================================================================
\section{Biblioteca de Runtime}
% =============================================================================

El compilador enlaza cada ejecutable con una biblioteca de runtime
(\texttt{hulk\_runtime.a}) que proporciona operaciones que son complejas de
implementar directamente en LLVM IR:

\begin{table}[H]
\centering
\begin{tabular}{ll}
\toprule
\textbf{Funcion} & \textbf{Descripcion} \\
\midrule
\texttt{hulk\_print(ptr)} & Imprime cadena + newline a stdout \\
\texttt{hulk\_str\_from\_number(f64)} & Convierte double a C string; enteros sin decimales \\
\texttt{hulk\_str\_concat(ptr, ptr)} & Concatenacion para el operador \texttt{@} \\
\texttt{hulk\_str\_concat\_space(ptr, ptr)} & Concatenacion con espacio para \texttt{@@} \\
\texttt{hulk\_str\_eq(ptr, ptr)} & Igualdad de cadenas por valor (no por puntero) \\
\texttt{hulk\_str\_size(ptr)} & Longitud en bytes de la cadena \\
\texttt{hulk\_vec\_alloc(n, sz)} & Alloca buffer: 8 bytes header + n*8 bytes elementos \\
\texttt{hulk\_vec\_get(ptr, i, sz)} & Puntero al elemento i con bounds check \\
\texttt{hulk\_vec\_size(ptr)} & Lee el header del vector y devuelve el count \\
\texttt{hulk\_range\_alloc(s, e)} & Crea iterador de rango $[s, e)$ \\
\texttt{hulk\_range\_next(ptr)} & Avanza el iterador; devuelve true si hay mas elementos \\
\texttt{hulk\_range\_current(ptr)} & Valor actual del iterador \\
\texttt{hulk\_rand()} & Numero aleatorio uniforme en $[0,1)$ \\
\texttt{hulk\_type\_error(ptr)} & Imprime error de tipo y aborta (para \texttt{as} fallido) \\
\bottomrule
\end{tabular}
\caption{Funciones de la biblioteca de runtime de HULK.}
\end{table}

\subsection{Layout de vectores en memoria}

Los vectores HULK se almacenan en memoria dinamica con la siguiente distribucion:

\begin{verbatim}
[ i64 count | elem[0] (8 bytes) | elem[1] (8 bytes) | ... | elem[n-1] ]
\end{verbatim}

Los primeros 8 bytes almacenan el numero de elementos como un entero de 64 bits.
Cada elemento ocupa 8 bytes independientemente de su tipo (numeros como
\texttt{double}, punteros como \texttt{i64} en plataformas 64-bit). Esto uniformiza
el acceso pero requiere que el codegen sepa el tipo de los elementos para emitir
el \texttt{load} correcto.

La funcion \texttt{hulk\_vec\_get} verifica los limites antes de retornar el puntero
al elemento:
\begin{lstlisting}[style=rust, caption={Bounds check en hulk\_vec\_get}]
pub extern "C" fn hulk_vec_get(
    vec: *mut u8, index: i32, _sz: i32
) -> *mut u8 {
    let size = unsafe { *(vec as *const i64) };
    if index < 0 || (index as i64) >= size {
        eprintln!("HULK error: indice {} fuera de rango [0, {})", index, size);
        std::process::abort();
    }
    unsafe { vec.add(8 + (index as usize) * 8) }
}
\end{lstlisting}

\subsection{Pipeline AOT completo}

Al invocar \texttt{./hulk source.hulk}:

\begin{enumerate}
  \item Analisis lexico + sintactico + semantico sobre el fuente.
  \item Generacion de IR LLVM para todas las declaraciones y el cuerpo principal.
  \item Adicion de un \texttt{main()} compatible con C que llama a
    \texttt{\_\_hulk\_entry()} y retorna 0.
  \item Optimizacion del IR con \texttt{mem2reg,instcombine,reassociate,simplifycfg}.
  \item Emision de objeto nativo a \texttt{/tmp/hulk\_program.o} via
    \texttt{TargetMachine::write\_to\_file}.
  \item Enlace con GCC: \texttt{gcc /tmp/hulk\_program.o hulk\_runtime.a -o ./output -lm}.
\end{enumerate}

El ejecutable resultante es un binario ELF x86-64 autocontenido que no requiere
LLVM en tiempo de ejecucion.

% =============================================================================
\section{Extension al Lenguaje: Vectores y Expresiones Generadoras}
\label{sec:extension}
% =============================================================================

La extension principal implementada es un \textbf{sistema de vectores con expresiones
generadoras}: construcciones sintacticas y semanticas que permiten crear y manipular
colecciones de valores de manera concisa, con soporte de primera clase en el sistema
de tipos.

\subsection{Motivacion}

HULK en su especificacion basica carece de estructuras de datos. Para escribir
algoritmos que operan sobre colecciones de valores, el programador se ve obligado a
usar tipos definidos por el usuario con logica de manejo de memoria manual. Esta
limitacion hace que tareas simples como calcular la media de N numeros o filtrar
valores sean verbosas y propensas a errores.

La extension de vectores cubre esta necesidad de forma minimalista y expresiva, siguiendo
la filosofia del lenguaje: todo es una expresion, y las operaciones comunes deben
poder escribirse en una sola linea.

\subsection{Sintaxis de la extension}

\subsubsection{Vectores explicitos}

Un vector se construye listando sus elementos entre corchetes, separados por comas:

\begin{lstlisting}[caption={Vector explicito con tipo inferido}]
let v = [1, 2, 3, 4, 5] in print(v[2]);   // imprime 3
let names = ["Alice", "Bob", "Carol"] in print(names[0]);
\end{lstlisting}

El tipo del vector se infiere del tipo de sus elementos. Si todos son \texttt{Number},
el tipo es \texttt{[Number]}. Si los tipos difieren pero tienen un ancestro comun, se
usa ese ancestro.

\subsubsection{Indexacion}

Los elementos se acceden con el operador \texttt{[i]}, donde \texttt{i} es una
expresion de tipo \texttt{Number}:

\begin{lstlisting}[caption={Indexacion con expresion}]
let v = [10, 20, 30] in {
    let i = 1 in v[i + 1]   // equivalente a v[2] = 30
}
\end{lstlisting}

El indice se convierte de \texttt{Number} (double) a entero en tiempo de compilacion
mediante \texttt{fptosi}. El runtime verifica que el indice sea no-negativo y menor
que el tamano del vector; en caso contrario, aborta con un mensaje descriptivo.

\subsubsection{Expresiones generadoras}

La sintaxis central de la extension es la \textbf{expresion generadora}, que crea un
nuevo vector evaluando una expresion para cada elemento de un iterable:

\begin{lstlisting}[caption={Expresiones generadoras}]
// Cuadrados de 0 a 4
let squares = [x^2 || x in range(0, 5)] in print(squares[3]); // 9

// Transformacion de otro vector
let v = [1, 2, 3] in
let doubled = [x * 2 || x in v] in
print(doubled[1]);  // 4
\end{lstlisting}

La sintaxis es \texttt{[body || var in iterable]}, donde:
\begin{itemize}
  \item \texttt{body} es cualquier expresion HULK (puede usar \texttt{var} y
    variables del entorno circundante).
  \item \texttt{var} es la variable de iteracion, local al cuerpo.
  \item \texttt{iterable} es un \texttt{Range} (del builtin \texttt{range(s,e)})
    o un \texttt{Vector} de cualquier tipo.
  \item El separador \texttt{||} es identico al operador OR logico pero en este
    contexto tiene semántica de generador (la gramatica distingue el contexto).
\end{itemize}

\subsubsection{Tamano y metodo \texttt{size}}

Tanto vectores como rangos exponen su tamano con la funcion \texttt{size}:

\begin{lstlisting}
let v = [3, 1, 4, 1, 5] in print(size(v));  // 5
\end{lstlisting}

\subsubsection{Tipo \texttt{[T]} en el sistema de tipos}

El tipo de un vector de \texttt{Number} se escribe \texttt{[Number]} y puede
aparecer en anotaciones de parametros y retornos:

\begin{lstlisting}[caption={Vector en firma de funcion}]
function sum(v: [Number]): Number =>
    let acc = 0 in {
        for (x in range(0, size(v))) acc := acc + v[x];
        acc
    };
\end{lstlisting}

La inferencia de tipos propaga el tipo del vector a traves de expresiones complejas.
El type checker verifica que la indexacion \texttt{v[i]} solo se aplique a valores
de tipo vector.

\subsection{Semantica de evaluacion: estrategia \textit{eager}}

Una decision fundamental en el diseno de colecciones es la \textbf{estrategia de
evaluacion}: \textit{eager} (inmediata) vs \textit{lazy} (diferida).

En la implementacion de HULK, los vectores y las expresiones generadoras se evaluan
\textbf{de forma eager}: todos los elementos se computan y almacenan en memoria en el
momento de la creacion del vector. Esto contrasta con otras aproximaciones:

\subsubsection{Comparacion con Python}

Python distingue entre listas (eager) y generadores (lazy). Las \textit{list
comprehensions} son eager:

\begin{lstlisting}[language=Python, caption={List comprehension en Python (eager)}]
squares = [x**2 for x in range(5)]  # crea lista inmediatamente
gen = (x**2 for x in range(5))      # generador lazy, no evalua hasta iteracion
\end{lstlisting}

La sintaxis de HULK \texttt{[body || var in iterable]} corresponde a las
\textit{list comprehensions} de Python (no a los generadores lazy), tanto en
comportamiento como en semantica de alcance: \texttt{var} es local al cuerpo.

La diferencia clave es que Python soporta condiciones (\texttt{if x > 0}) y
multiples iteradores (\texttt{for x in xs for y in ys}), mientras que HULK limita
la expresion generadora a un solo iterador sin filtro. Esta simplificacion fue
intencional: mantiene la gramatica LALR(1) sin ambiguedades y reduce la complejidad
del codegen.

\subsubsection{Comparacion con Haskell}

Haskell tiene \textit{list comprehensions} con sintaxis similar a la notacion
matematica de conjuntos:

\begin{lstlisting}[language=Haskell, caption={List comprehension en Haskell (lazy)}]
squares = [x^2 | x <- [0..4]]    -- notacion matematica: {x^2 | x in {0..4}}
evens   = [x | x <- [1..100], even x]  -- con filtro
\end{lstlisting}

La evaluacion en Haskell es \textit{lazy por defecto}: los elementos de la lista se
computan solo cuando son demandados. Esto permite listas infinitas:
\texttt{nats = [0..]} es valido en Haskell porque los elementos se computan
incrementalmente.

HULK usa evaluacion eager, lo que hace que el tamano del vector deba conocerse antes
de la iteracion. Para rangos (\texttt{range(s,e)}), el count se calcula como
\texttt{floor(e - s)}; para vectores, se consulta el header. Esta decision es adecuada
para un lenguaje compilado con modelo de memoria simple: la evaluacion lazy requiere
thunks (cierres de suspension) que anadiran complejidad al modelo de objetos.

La notacion \texttt{|| var in iterable} en HULK es deliberadamente similar a la
notacion de Haskell \texttt{| var <- list}, evocando la notacion matematica
de construccion de conjuntos $\{body \mid var \in S\}$.

\subsubsection{Comparacion con Scala}

Scala ofrece comprensiones via \texttt{for...yield}:

\begin{lstlisting}[language=Scala, caption={For-comprehension en Scala}]
val squares = for (x <- 0 until 5) yield x * x  // Vector[Int]
val pairs   = for {
  x <- 1 to 3
  y <- 1 to x
} yield (x, y)  // multiples generadores
\end{lstlisting}

Scala es eager para colecciones concretas (\texttt{Vector}, \texttt{List}) y lazy
para vistas (\texttt{.view}). La sintaxis \texttt{for...yield} es mas verbosa pero
mas flexible: soporta multiples generadores, guardas y definiciones locales. HULK
simplifica esto a un solo generador, lo que se traduce en un patron de acceso
frecuente (transformar una coleccion) sin la verbosidad de Scala.

\subsubsection{Comparacion con C++}

C++20 introduce \textit{ranges} y vistas lazy:

\begin{lstlisting}[language=C++, caption={Ranges y views en C++20 (lazy)}]
auto squares = std::views::iota(0, 5)
             | std::views::transform([](int x){ return x*x; });
// squares es lazy: no evalua hasta iteracion
std::vector<int> v(squares.begin(), squares.end()); // materializa
\end{lstlisting}

C++ favorece la evaluacion lazy con materializacion explicita, optimizada para
evitar allocaciones intermedias. Para HULK, que carece de iteradores genericos y
de un sistema de tipos con lifetimes, la aproximacion C++ seria demasiado compleja.

\subsubsection{Resumen comparativo}

\begin{table}[H]
\centering
\begin{tabular}{lllll}
\toprule
\textbf{Lenguaje} & \textbf{Sintaxis} & \textbf{Evaluacion} & \textbf{Filtros}
  & \textbf{Tipos} \\
\midrule
HULK    & \texttt{[e || x in it]} & Eager  & No  & Estatico inferido \\
Python  & \texttt{[e for x in it if c]} & Eager/Lazy & Si & Dinamico \\
Haskell & \texttt{[e | x <- xs, c]} & Lazy & Si & Estatico inferido \\
Scala   & \texttt{for x <- xs yield e} & Eager/Lazy & Si & Estatico inferido \\
C++20   & \texttt{xs | transform(f)} & Lazy & Si & Estatico explicito \\
\bottomrule
\end{tabular}
\caption{Comparativa de expresiones generadoras entre lenguajes.}
\end{table}

\subsection{Implementacion del codegen para vectores}

\subsubsection{Vector explicito}

La generacion de codigo para un vector explicito \texttt{[e0, e1, ..., en-1]}:

\begin{enumerate}
  \item Llama a \texttt{hulk\_vec\_alloc(n, 8)} para reservar el buffer.
  \item Para cada elemento $i$: evalua $e_i$, llama a \texttt{hulk\_vec\_get} para
    obtener el puntero al slot $i$, y almacena el valor.
  \item Retorna el puntero al buffer como \texttt{CgValue::Vector}.
\end{enumerate}

\subsubsection{Expresion generadora}

La generacion para \texttt{[body || var in iterable]} es mas compleja:

\begin{enumerate}
  \item Evalua \texttt{iterable} para obtener el puntero al rango o vector fuente.
  \item Calcula el count: para rangos, \texttt{floor(end - start)} con clamp a
    $[0, \texttt{i32::MAX}]$ para evitar \textit{LLVM poison} por fptosi fuera
    de rango; para vectores, llama a \texttt{hulk\_vec\_size}.
  \item Llama a \texttt{hulk\_vec\_alloc(count, 8)} para el resultado.
  \item Genera un bucle LLVM con variable de indice:
    \begin{itemize}
      \item \texttt{gen\_cond}: carga el indice, compara con count, salta a
        \texttt{gen\_body} o \texttt{gen\_end}.
      \item \texttt{gen\_body}: carga o calcula el elemento fuente, lo asigna a
        \texttt{var}, evalua \texttt{body}, almacena el resultado en el slot
        correspondiente del vector resultado, incrementa el indice.
      \item \texttt{gen\_end}: retorna el puntero al vector resultado.
    \end{itemize}
\end{enumerate}

El clamp del count para rangos merece especial atencion: la instruccion \texttt{fptosi}
de LLVM produce \textit{poison value} si el operando esta fuera del rango del tipo
entero destino. Un rango con \texttt{start > end} daria \texttt{end - start < 0},
lo que convertiria a un entero negativo enorme. El codigo inserta:

\begin{lstlisting}[style=llvmir, caption={Clamp del count en generador de rango}]
%diff   = fsub double %end, %start
; si diff < 0 (start > end), forzar a 0.0
%is_neg = fcmp ult double %diff, 0.0
%nn     = select i1 %is_neg, double 0.0, double %diff
; si diff > i32max, forzar a i32max (evita poison en fptosi)
%hi     = fcmp ogt double %nn, 2147483647.0
%clamp  = select i1 %hi, double 2147483647.0, double %nn
%count  = fptosi double %clamp to i32
\end{lstlisting}

\subsection{El tipo Range como iterable}

Ademas de los vectores, la extension incluye el tipo \texttt{Range} como un iterable
de primera clase. El \texttt{Range} es un objeto en heap con el layout:

\begin{verbatim}
[ f64 start | f64 end | f64 current ]
\end{verbatim}

\texttt{current} arranca en \texttt{start - 1.0}. Cada llamada a
\texttt{hulk\_range\_next} incrementa \texttt{current} en 1.0 y retorna
\texttt{current < end}. Cuando la condicion falla, la iteracion termina.

Los bucles \texttt{for (x in range(s, e))} se compilan a un patron de iteracion
sobre \texttt{Range}:

\begin{lstlisting}[style=llvmir, caption={Compilacion de for (x in range(0, 3)) body}]
  %range  = call ptr @hulk_range_alloc(double 0.0, double 3.0)
loop_cond:
  %cont   = call i1 @hulk_range_next(ptr %range)
  br i1 %cont, label %loop_body, label %loop_end
loop_body:
  %x_val  = call double @hulk_range_current(ptr %range)
  ; ... body usando %x_val ...
  br label %loop_cond
loop_end:
\end{lstlisting}

\subsection{Justificacion de las decisiones de diseno}

\subsubsection{Por que eager y no lazy?}

La evaluacion eager se eligio por tres razones:

\begin{enumerate}
  \item \textbf{Simplicidad del modelo de memoria}: con evaluacion eager, el tamano
    del vector se conoce antes de la asignacion. Esto permite una sola llamada a
    \texttt{malloc} por vector. La evaluacion lazy requeriria thunks o continuaciones,
    estructuras de datos adicionales y un protocolo de evaluacion mas complejo.
  \item \textbf{Rendimiento predecible}: los compiladores de lenguajes lazy como GHC
    invierten considerable esfuerzo en evitar los overheads de la thunkification.
    Para un compilador educativo que ya usa LLVM, eager produce IR mas simple y
    mas facil de optimizar.
  \item \textbf{Consistencia con el modelo imperativo}: HULK tiene efectos de lado
    (print, asignacion destructiva). La evaluacion lazy en presencia de efectos produce
    comportamientos sorprendentes (el orden de evaluacion no es el orden textual).
    Eager garantiza que los efectos se producen en orden de izquierda a derecha.
\end{enumerate}

\subsubsection{Por que \texttt{||} como separador?}

La notacion \texttt{[body || var in iterable]} fue disenada para evocar la notacion
matematica de comprension de conjuntos $\{f(x) \mid x \in S\}$ y la sintaxis de
Haskell \texttt{[f x | x <- xs]}. El simbolo \texttt{||} es el mismo que el OR logico,
pero la gramatica LALR lo distingue por contexto (solo es un separador dentro de
\texttt{[...]}). Alternativas consideradas:

\begin{itemize}
  \item \texttt{[body for var in iterable]}: mas legible para programadores Python,
    pero introduce un nuevo keyword \texttt{for} en un contexto diferente al de los
    bucles, lo que crea ambiguedad en la gramatica LALR.
  \item \texttt{[body : var in iterable]}: colisiona con la notacion de tipo
    \texttt{x : T} ya existente en la gramatica.
  \item \texttt{[body | var in iterable]}: un solo \texttt{|} colisiona con las
    alternativas de la gramatica en la notacion BNF (confunde al desarrollador al
    leer el archivo de gramatica).
\end{itemize}

\subsubsection{Por que vectores de tamano fijo?}

Los vectores HULK no tienen \texttt{push} ni \texttt{pop}: una vez creados, su tamano
no cambia. Esta decision se tomo para:

\begin{enumerate}
  \item \textbf{Mantener el modelo de memoria simple}: un buffer fijo con header de
    count es trivial de implementar y de razonar. Vectores dinamicos requieren
    estrategias de redimension (duplicar capacidad, etc.) que duplican la complejidad
    del runtime.
  \item \textbf{Conservar la semantica de expresion}: en HULK todo es una expresion.
    Un vector dinamico con mutacion es un objeto con estado, lo que rompe la
    composicion funcional que permite escribir \texttt{print([x\^{}2 || x in range(0,5)][3])}.
\end{enumerate}

% =============================================================================
\section{Cobertura de la Especificacion HULK}
% =============================================================================

La siguiente tabla muestra el estado de implementacion de cada seccion de la
especificacion oficial de HULK:

\begin{table}[H]
\centering
\begin{tabular}{clcc}
\toprule
\textbf{Seccion} & \textbf{Feature} & \textbf{Req.} & \textbf{Estado} \\
\midrule
16.1 & Expresiones aritmeticas y literales & Si & Completo \\
16.2 & Variables y \texttt{let-in} & Si & Completo \\
16.3 & Funciones y llamadas & Si & Completo \\
16.4 & Condicionales \texttt{if-elif-else} & Si & Completo \\
16.5 & Ciclos \texttt{while} & Si & Completo \\
16.6 & Tipos, herencia, metodos, \texttt{self} & Si & Completo \\
16.7 & Protocolos y conformidad estructural & Si & Completo \\
16.8 & Operadores \texttt{is} y \texttt{as} & Si & Completo \\
16.9 & Vectores y colecciones & No & \textbf{Extension} \\
16.10 & Expresiones generadoras & No & \textbf{Extension} \\
16.11 & Iteradores de rango & No & \textbf{Extension} \\
--- & Operador \texttt{base()} & No & Implementado \\
--- & Escape sequences en strings & Si & Implementado \\
--- & Constantes \texttt{PI} y \texttt{E} & Si & Implementado \\
--- & Funciones matematicas builtin & Si & Implementado \\
\midrule
& \textbf{Total obligatorio cubierto} & & \textbf{100\%} \\
\bottomrule
\end{tabular}
\caption{Cobertura de la especificacion HULK por seccion.}
\end{table}

El compilador cubre el \textbf{100\%} de los features obligatorios. Las extensiones
(secciones~16.9--16.11 y funcionalidades adicionales) amplian el lenguaje mas alla
de los requisitos minimos.

\subsection{Ejemplos de programas HULK completos}

\subsubsection{Herencia y polimorfismo}

\begin{lstlisting}[caption={Herencia y dispatch virtual}]
type Animal(name: String) {
    name = name;
    speak(): String => "...";
    greet(): String => "I am " @@ self.name @@ " and I say: " @@ self.speak();
}

type Dog(name: String) inherits Animal(name) {
    speak(): String => "Woof!";
}

type Cat(name: String) inherits Animal(name) {
    speak(): String => "Meow!";
}

let animals = [new Dog("Rex"), new Cat("Whiskers")] in {
    for (a in animals) print(a.greet());
}
\end{lstlisting}

\subsubsection{Protocolos}

\begin{lstlisting}[caption={Protocolo Printable}]
protocol Printable {
    to_str(): String;
}

type Point(x: Number, y: Number) {
    x = x; y = y;
    to_str(): String => "(" @ x @ ", " @ y @ ")";
}

function show(item: Printable) => print(item.to_str());

show(new Point(3, 4));
\end{lstlisting}

\subsubsection{Extension: generadores}

\begin{lstlisting}[caption={Fibonacci con vectores}]
function fibonacci(n: Number): [Number] =>
    let fibs = [0, 1] in
    [
        if (i < 2) fibs[i]
        else fibs[i-1] + fibs[i-2]
        || i in range(0, n)
    ];

let f = fibonacci(10) in
for (x in range(0, size(f))) print(f[x]);
\end{lstlisting}

% =============================================================================
\section{Pruebas y Evaluacion}
% =============================================================================

\subsection{Suite de pruebas automatizadas}

El compilador incluye pruebas automatizadas ejecutables con \texttt{cargo test}:

\begin{table}[H]
\centering
\begin{tabular}{lrl}
\toprule
\textbf{Modulo} & \textbf{Tests} & \textbf{Cobertura} \\
\midrule
Lexer    & 87  & Todos los tipos de token, edge cases, secuencias de escape \\
Parser   & 112 & Todos los no-terminales de la gramatica, precedencia \\
Semantico & 98 & Inferencia de tipos, protocolos, errores de tipo \\
Codegen JIT & 87 & Aritmetica, flujo, OOP, vectores, rangos, builtins \\
\midrule
\textbf{Total} & \textbf{384} & \\
\bottomrule
\end{tabular}
\caption{Distribucion de las pruebas automatizadas.}
\end{table}

\subsection{Suite oficial del profesor}

La suite de pruebas oficial esta organizada en casos que verifican la salida estandar
exacta del programa. El script \texttt{run\_tests.sh} compila cada fichero
\texttt{tests/professor\_suite/test\_XX\_nombre.hulk} y compara la salida con el
valor esperado. Los tests cubren:

\begin{itemize}
  \item Expresiones aritmeticas y precedencia.
  \item Cadenas y operadores \texttt{@}, \texttt{@@}.
  \item Funciones recursivas e iterativas.
  \item Herencia, override, despacho virtual.
  \item Protocolos con multiples tipos conformantes.
  \item Operadores \texttt{is} y \texttt{as}.
  \item Vectores y expresiones generadoras.
  \item Casos de error con formato de salida requerido.
\end{itemize}

\subsection{Contrato de errores}

El compilador sigue el contrato de interfaz del profesor: los errores se reportan
por \textbf{stderr} con el prefijo \texttt{!\ } y el formato
\texttt{(linea,col) FASE: mensaje}. El proceso termina con codigo de salida $\neq 0$
cuando hay errores. Esto permite al CI verificar automaticamente los casos de error.

% =============================================================================
\section{Decisiones de Diseno}
\label{sec:decisions}
% =============================================================================

\subsection{LALR(1) artesanal vs generador externo}

Se eligio implementar el parser desde cero (gramatica como estructura de datos,
tablas construidas en Rust) en lugar de usar YACC, LALRPOP o ANTLR por:

\begin{enumerate}
  \item \textbf{Control total sobre la gramatica}: es posible inspeccionarla,
    imprimirla y depurarla en tiempo de ejecucion. Con un generador externo, los
    conflictos se reportan sin contexto de que producciones los causan.
  \item \textbf{Sin dependencias de build externas}: la gramatica se compila con
    el proyecto, sin requerir una herramienta separada.
  \item \textbf{Valor educativo}: implementar LALR(1) completo es el objetivo central
    de la asignatura de Compiladores. Usar un generador elude el aprendizaje.
\end{enumerate}

El costo de esta decision es una implementacion mas larga ($\sim$800 LOC solo para el
modulo LALR) y un tiempo de startup ligeramente mayor (las tablas se construyen en
cada invocacion, aunque esto podria cachearse en disco).

\subsection{Backend LLVM vs interprete vs bytecode propio}

Se eligio LLVM por:

\begin{itemize}
  \item \textbf{Codigo nativo de calidad}: LLVM produce codigo comparable al de GCC
    para estructuras comunes. Un interprete o VM propio seria varias ordenes de magnitud
    mas lento.
  \item \textbf{Portabilidad}: el mismo IR de LLVM puede compilarse para x86-64,
    ARM, WebAssembly u otras plataformas simplemente cambiando el TargetMachine.
  \item \textbf{Optimizaciones gratuitas}: el pipeline de optimizacion de LLVM
    (\texttt{mem2reg}, constant propagation, etc.) mejora sustancialmente el codigo
    generado sin esfuerzo adicional.
\end{itemize}

\subsection{Vtables vs tagged union vs boxed types}

Para implementar el polimorfismo se consideraron tres alternativas:

\begin{enumerate}
  \item \textbf{Tagged union}: cada objeto es una union discriminada con un tag de tipo.
    Simple pero ineficiente: el tamano de la union debe ser el del tipo mas grande.
  \item \textbf{Boxed types}: cada objeto es un puntero a un area de heap con layout
    dinamico. Requiere un sistema de tipos en runtime mas complejo.
  \item \textbf{Vtables (elegida)}: cada objeto tiene un puntero a una tabla de
    funciones. Es la tecnica usada por C++, Java (con JIT), y Rust (para trait objects).
    El costo por llamada es una carga adicional de puntero (tipicamente en cache);
    el beneficio es que el layout de cada tipo es estatico y conocido.
\end{enumerate}

\subsection{Enlace dinamico de LLVM vs estatico}

El compilador enlaza \texttt{libLLVM-17.so} dinamicamente. El enlace estatico produciria
ejecutables de 50--200 MB y seria mas lento de enlazar. El dinamico mantiene el
binario en pocos MB y comparte la biblioteca con otras herramientas del sistema.

% =============================================================================
\section{Limitaciones Conocidas y Trabajo Futuro}
% =============================================================================

\subsection{Limitaciones actuales}

\begin{enumerate}
  \item \textbf{Funciones anonimas (lambdas)}: HULK soporta funciones de primera clase
    en su especificacion (seccion~16.12). La implementacion actual solo soporta funciones
    con nombre. La adicion de lambdas requeriria \textit{closures} con captura del
    entorno, lo que implica un cambio en el modelo de objetos.
  \item \textbf{Gestion de memoria}: el compilador usa una estrategia de solo-alloc:
    los strings y objetos se asignan en el heap y nunca se liberan. Para programas de
    larga ejecucion con muchas operaciones de cadena esto provoca crecimiento ilimitado
    de memoria. Una mejora futura podria ser un \textit{arena allocator} simple o
    un contador de referencias.
  \item \textbf{Recuperacion de errores en el parser}: el parser se detiene al primer
    error sintactico. Un parser con recuperacion podria reiniciar en el siguiente
    punto de sincronizacion (\texttt{;} o \texttt{\}}) para reportar multiples errores.
  \item \textbf{Mensajes de error en el codegen sin span}: los errores de generacion
    de codigo no incluyen la posicion en el fuente. Las fases anteriores si incluyen
    spans; propagar estos al codegen mejoraria la experiencia de usuario.
  \item \textbf{Vectores sin mutacion estructural}: los vectores son de tamano fijo.
    Agregar \texttt{push}/\texttt{pop} requeriria un buffer dinamico en el runtime
    similar al de \texttt{Vec<T>} de Rust.
\end{enumerate}

\subsection{Trabajo futuro}

\begin{itemize}
  \item \textbf{Closures y lambdas}: representar funciones de primera clase con
    un puntero a funcion + puntero a entorno capturado.
  \item \textbf{Generadores lazy}: evaluar las expresiones generadoras bajo demanda,
    permitiendo trabajar con colecciones conceptualmente infinitas.
  \item \textbf{Inferencia de tipos full Hindley-Milner}: el sistema actual infiere
    tipos a traves de expresiones simples pero no resuelve constraints polimorficos
    complejos. H-M permitiria funciones genericas como
    \texttt{function map(v: [A], f: A -> B): [B]}.
  \item \textbf{Optimizaciones especificas de HULK}: inlining de metodos monomorfonicos,
    escape analysis para objetos de corta vida (alloca en stack en lugar de heap).
  \item \textbf{Soporte para WASM}: cambiar el TargetMachine de LLVM para emitir
    WebAssembly permitiria ejecutar programas HULK en el navegador.
\end{itemize}

% =============================================================================
\section{Conclusiones}
% =============================================================================

Se implemento un compilador completo para el lenguaje HULK en Rust con backend LLVM~17.
El compilador cubre el 100\% de los features obligatorios de la especificacion
(secciones~1--16.8) y extiende el lenguaje con un sistema de vectores y expresiones
generadoras que modifica tanto la sintaxis como la semantica del lenguaje.

Las decisiones tecnicas mas significativas --- LALR(1) artesanal, vtables con type
tags DFS, evaluacion eager para generadores --- estan justificadas por compromisos
de simplicidad de implementacion, rendimiento predecible, y coherencia con el modelo
de ejecucion de HULK.

La extension de vectores con expresiones generadoras coloca a HULK en una familia de
lenguajes con soporte de primera clase para transformaciones de colecciones, junto a
Python, Haskell y Scala, con la particularidad de integrarse en un lenguaje tipado
estaticamente con inferencia de tipos y compilacion a codigo nativo.

El resultado es un compilador que produce ejecutables nativos x86-64 sin dependencias
en tiempo de ejecucion (salvo \texttt{libc} y la biblioteca de runtime propia), con
rendimiento comparable al codigo C equivalente para programas con aritmetica y OOP.

% =============================================================================
\begin{thebibliography}{99}
% =============================================================================

\bibitem{aho2006}
A. V. Aho, M. S. Lam, R. Sethi, J. D. Ullman.
\textit{Compilers: Principles, Techniques, and Tools} (2nd ed.).
Addison-Wesley, 2006. [Dragon Book]

\bibitem{thompson1968}
K. L. Thompson.
``Programming Techniques: Regular Expression Search Algorithm.''
\textit{Communications of the ACM}, 11(6):419--422, 1968.

\bibitem{deremer1969}
F. L. DeRemer.
\textit{Practical Translators for LR(k) Languages}.
PhD thesis, MIT, 1969.

\bibitem{lattner2004}
C. Lattner, V. Adve.
``LLVM: A Compilation Framework for Lifelong Program Analysis and Transformation.''
\textit{Proceedings of CGO}, 2004.

\bibitem{bjarne1994}
B. Stroustrup.
\textit{The Design and Evolution of C++}.
Addison-Wesley, 1994.

\bibitem{inkwell}
inkwell: Safe Rust bindings for LLVM.
\url{https://github.com/TheDan64/inkwell}. Accedido: 2025.

\bibitem{pierce2002}
B. C. Pierce.
\textit{Types and Programming Languages}.
MIT Press, 2002.

\bibitem{hudak1989}
P. Hudak.
``Conception, Evolution, and Application of Functional Programming Languages.''
\textit{ACM Computing Surveys}, 21(3):359--411, 1989.

\bibitem{peytonjones2003}
S. Peyton Jones (ed.).
\textit{Haskell 98 Language and Libraries: The Revised Report}.
Cambridge University Press, 2003.

\bibitem{odersky2008}
M. Odersky, L. Spoon, B. Venners.
\textit{Programming in Scala}.
Artima Press, 2008.

\bibitem{klabnik2019}
S. Klabnik, C. Nichols.
\textit{The Rust Programming Language}.
No Starch Press, 2019.

\bibitem{matcom2024}
Equipo MatCom.
\textit{HULK Language Specification}.
Facultad de Matematica y Computacion, Universidad de La Habana, 2024.
\url{https://matcom.github.io/hulk/}

\end{thebibliography}

\end{document}
